\subsection{The normalized ordinary cubic Massey condition}
\label{sec:kt-cubic-massey-condition}

The next condition is a concrete higher lifting condition, not a
vanishing theorem for the original Selmer group.  Its rational
class-group character is different from the character controlling
that original negative root.  The local matching conditions are
retained throughout; raw global $H^2$ is not substituted for the
ordinary obstruction space.

Let $p>3$, let $e\ge2$ be even with $(p-1)\nmid2e$, let
$D<0$ be fundamental and prime to $p$, and use
\[
 \eta=\chi_D,\quad \psi=\omega^e,\quad
 \alpha=\eta\psi,\quad\beta=\psi^{-1},\quad
 S=\{\infty,p\}\cup\{q:q\mid D\}.
                                                               \tag{CM.1}
\]
Let $k/\mathbf F_p$ be finite.  Assume the two auxiliary units
$B_{1,\eta}$ and $B_{1,\eta\omega^e}$.
The third unit at exponent $2e$ is not needed at this degree.
Fix continuous cocycles
\[
 u\in Z^1(G_{\mathbf Q,S},k(\alpha)),\qquad
 v\in Z^1(G_{\mathbf Q,S},k(\beta)),
                                                               \tag{CM.2}
\]
where $[v]$ is zero at $p$ and is unramified on all other finite
inertia groups.  These are the two directions in
KT-GLOBAL-QUADRATIC-MIXED-PERSISTENCE.

\paragraph{The specified graded equivariance.}
Put
\[
\begin{gathered}
 X=E_{12}+E_{23},\quad Y=E_{21}-E_{32},\quad
 Z=\operatorname{diag}(1,-2,1),\\
 U=E_{12}-E_{23},\qquad V=E_{21}+E_{32}.
\end{gathered}                                                 \tag{CM.3}
\]
For $T(a)=\operatorname{diag}(a,1,a^{-1})$, require the algebraic
$\mathbf G_m$-equivariance that sends $(s,t)$ to $(as,a^{-1}t)$
and acts on matrices by conjugation by $T(a)$.  This is an
integer grading, not merely equivariance for the finite group
$k^\times$.  Also require equivariance under
\[
 \kappa(g)=D_c\theta(g)D_c,\qquad
 D_c=\operatorname{diag}(-1,1,-1),\qquad
 \kappa(s,t)=(s,-t),                                    \tag{CM.4}
\]
where $\theta(g)=J(g^t)^{-1}J$ is the unitary outer automorphism.
It fixes the residual representation $r_0$.
On the five Lie vectors in (CM.3), the signs of $\kappa$ are
$+,-,-,-,+$.  Their Galois characters are respectively
$\alpha,\beta,\eta,\psi,\eta\psi^{-1}$.

Use $A_2=k[s,t]/(s,t)^3$ and $A_3=k[s,t]/(s,t)^4$.
The determinant-one, graded, $\kappa$-equivariant logarithms
through degree three necessarily have the form
\[
\begin{split}
 L(g)={}&s u(g)X+t v(g)Y+st w(g)Z\\
        &+s^2t\,m(g)U+st^2\,n(g)V.
\end{split}                                                   \tag{CM.5}
\]
Indeed the weight-two Lie line has $\kappa$ sign minus,
whereas $s^2$ has sign plus, and similarly for $t^2$.
At bidegree $(1,1)$ the trace-zero diagonal space of sign minus
is exactly $kZ$.  At degrees $(2,1)$ and $(1,2)$ the permitted
lines are $kU$ and $kV$.  There are no weight-three Lie lines.
Thus (CM.5) specifies the entire equivariant lifting problem,
rather than omitting other permitted pure-square corrections.

\begin{theorem}[KT-ORDINARY-CUBIC-MASSEY-CONDITION]
For the data (CM.1)--(CM.2), choose the ordinary quadratic
defining cochain $w$ supplied by
KT-GLOBAL-QUADRATIC-MIXED-PERSISTENCE.  There is a well-defined
ordinary cubic obstruction
\[
                  \mathfrak o_3(u,v)\in\mathcal O_\psi.          \tag{CM.6}
\]
Here $\mathcal O_\psi$ is the relative cochain obstruction space
defined below.  Put $F=\mathbf Q(\mu_p)$,
$\Delta_0=\operatorname{Gal}(F/\mathbf Q)$,
$A_F=\operatorname{Cl}(F)\{p\}$, and $\delta=\omega^{1-e}$.
Then
\[
 \dim_k\mathcal O_\psi
   =\dim_k\operatorname{Hom}_{\Delta_0}
                         (A_F/pA_F,k(\delta)).                 \tag{CM.7}
\]
The class (CM.6) is zero if and only if the given quadratic
family lifts over $A_3$ with the grading and involution
(CM.4), the exact ordered ordinary condition at $p$, and
unramified restriction to $G_E$ away from $p$.

Its image in raw global $H^2(k(\psi))$ is
$-[\Omega_+]$, where $\Omega_+$ is the explicit cochain
(CM.10) below.  It lies in the locally trivial subgroup.
If $b\in\{2,4,\ldots,p-3\}$ is congruent to $e$ modulo $p-1$
and $B_b/b$ is a $p$-adic unit, then
$\mathcal O_\psi=0$: every such quadratic family lifts to
this cubic order.  This extra rational unit removes the cubic
obstruction; it does not force $[v]=0$.
\end{theorem}

\begin{proof}
\paragraph{The quadratic defining cochain has no ordinary indeterminacy.}
The equation for $w$ is
\[
 w(gh)-w(g)-\eta(g)w(h)
  =\tfrac12\{u(g)\beta(g)v(h)-v(g)\alpha(g)u(h)\}.         \tag{CM.8}
\]
For a fixed local primitive $v(g)=\beta(g)d-d$ at $p$,
the ordinary condition requires the class of
$w_p=w-d(1+\beta)u/2$ to be unramified.
Two choices of $w$ satisfying these conditions differ by an
$\eta$-cocycle whose localization is unramified.  The
isomorphism (MC.3) makes this cocycle a global coboundary.
Conjugation by $\exp(st aZ)$ realizes that coboundary and
preserves all the indicated local conditions and equivariances.
Consequently the quadratic family is unique up to this
equivalence, once $u,v$ are fixed.  This is stronger than
merely saying that an unrestricted Massey product is defined.

\paragraph{The exact cubic equations.}
For a pair $g,h$, abbreviate
\[
\begin{array}{lll}
 A=u(g),&B=v(g),&C=w(g),\\
 a=\alpha(g)u(h),&b=\beta(g)v(h),&c=\eta(g)w(h),
\end{array}
\qquad F_{g,h}=Ab-Ba.                                   \tag{CM.9}
\]
Since
$[X,Y]=Z$, $[X,Z]=-3U$ and $[Y,Z]=3V$, the degree-three
Baker--Campbell--Hausdorff formula gives
\[
\begin{split}
 \Omega_+(g,h)&=-\tfrac32(Ac-Ca)
                         -\tfrac14(A-a)F_{g,h},\\
 \Omega_-(g,h)&=\tfrac32(Bc-Cb)
                         +\tfrac14(B-b)F_{g,h}.
\end{split}                                                   \tag{CM.10}
\]
The equations for the two remaining cochains are exactly
\[
\begin{split}
 m(gh)-m(g)-\psi(g)m(h)&=\Omega_+(g,h),\\
 n(gh)-n(g)-\eta(g)\psi(g)^{-1}n(h)&=\Omega_-(g,h).
\end{split}                                                   \tag{CM.11}
\]
The two right sides are continuous two-cocycles in the displayed
characters.  This follows either by substitution using (CM.8)
or by associativity: the degree-three kernel is an abelian
module, and the defect of a lift through a square-zero kernel
is a two-cocycle.  All exponential and BCH expressions here
are finite polynomial identities; the denominators divide
$12$, which is invertible.

\paragraph{The twisted cubic direction is removable.}
KT-CUBIC-TWISTED-COHOMOLOGY gives
\[
\begin{gathered}
 H^2(G_{\mathbf Q,S},k(\eta\psi^{-1}))=0,\\
 H^1(G_{\mathbf Q,S},k(\eta\psi^{-1}))
       \xrightarrow{\ \sim\ }H^1(\mathbf Q_p,k(\eta\psi^{-1})).
\end{gathered}                                                 \tag{CM.12}
\]
It also gives zero local $H^1$ in that character at every
$q\mid D$.  Thus the equation for $n$ has a global solution.
After the same local negative-root conjugation used at degree
two, its difference from an ordinary local lift at $p$ is a
cocycle in $k(\eta\psi^{-1})$.  Correct it by a global cocycle
using (CM.12).  At $q\mid D$ any difference from an
unramified-over-$E$ local lift is a coboundary.  Hence this
direction imposes no further global or local condition.
The two cubic directions are independent at this truncation,
because their interactions have degree at least four.

\paragraph{Local lifts of the rational direction exist.}
At $p$, $H^2(\mathbf Q_p,k(\psi))=0$, since its Cartier
dual has nontrivial inertia character $\omega^{1-e}$.
After removing the lower first-order term, the quadratic
family is upper triangular with fixed inertial diagonal.
Its cubic upper correction therefore exists, with arbitrary
local $H^1(k(\psi))$ difference between choices.

At $q\mid D$, local $H^1$ in both $\alpha$ and $\eta$ is
zero.  The degree-two argument therefore conjugates the
family to $\exp(t v(g)Y)r_0(g)$, with $v$ unramified.
This last formula is itself a homomorphism to every
truncation, because $Y$ spans an abelian Lie line and $v$
is a cocycle.  Extending the finite conjugating exponentials
gives an acceptable cubic local lift.  The character $\psi$
is unramified at $q$, and the difference between two such
local lifts, modulo local conjugation, is precisely an
unramified $H^1(k(\psi))$ class.  Indeed after the inertia
of $E$ is killed, the remaining inertia quotient has order
at most two; its positive-degree $p$-cohomology is zero.
There is no ramification outside $S$.

\paragraph{The relative obstruction space.}
Let $M=k(\psi)$, and set
\[
 L_p=H^1(\mathbf Q_p,M),\qquad
 L_q=H^1_{\rm unr}(\mathbf Q_q,M)\quad(q\mid D),
 \qquad L_\infty=0.                                    \tag{CM.13}
\]
Use the ordinary cochain differential
$df(g,h)=f(g)+g f(h)-f(gh)$.
Define $\mathcal O_\psi$ as the vector space of pairs
\[
 (c,(z_v)_{v\in S}),\qquad
 c\in Z^2(G_{\mathbf Q,S},M),\qquad dz_v=c|_{G_{\mathbf Q_v}},
                                                               \tag{CM.14}
\]
modulo pairs
$(df,(f|_{G_{\mathbf Q_v}}+\ell_v)_v)$, where $f$ is a
global continuous one-cochain and each $\ell_v$ is a local
one-cocycle with class in $L_v$.  Local coboundaries are
included in this equivalence.  Its defining exact sequence is
\[
\begin{split}
0\longrightarrow&
 \operatorname{coker}\!\left(
 H^1(G_{\mathbf Q,S},M)\longrightarrow
       \bigoplus_{v\in S}H^1(\mathbf Q_v,M)/L_v\right)\\
\longrightarrow&\mathcal O_\psi
\longrightarrow\Sha^2_S(\mathbf Q,M)\longrightarrow0 .
\end{split}                                                   \tag{CM.15}
\]
This follows directly: forgetting the local primitives gives
the last map; for an exact global cocycle, subtract a global
primitive and retain the local one-cocycle classes.

Take $c=-\Omega_+$ and choose the local primitives supplied
by the acceptable local cubic lifts.  Their classes modulo
the allowed local differences do not depend on the choices.
The resulting pair defines $\mathfrak o_3(u,v)$.
It vanishes precisely when a global $m$ solving (CM.11)
can match all those local lifts modulo (CM.13).
Together with the already removable $n$ direction, this
proves the asserted necessary and sufficient condition.
Changing $w$ by its permitted global coboundary conjugates
the lifting problem and acts trivially on the central
degree-three kernel, so does not change the obstruction.
Changing cocycle representatives of $u,v$ is likewise
realized by the first-order conjugations and does not
change its equivalence class.

\paragraph{Its exact rational arithmetic dimension.}
The Cartier dual is $M^D=k(\delta)$.
At $q\mid D$, local unramified orthogonality gives
$L_q^\perp=H^1_{\rm unr}(\mathbf Q_q,M^D)$.
At $p$ it gives $L_p^\perp=0$.
Put
\[
 Q=\ker\!\left(H^1(G_{\mathbf Q,S},M^D)
       \longrightarrow\bigoplus_{v\in S}
                   H^1(\mathbf Q_v,M^D)/L_v^\perp\right).
                                                               \tag{CM.16}
\]
The Poitou--Tate exact sequence and the perfect local
pairings give
\[
 \dim\operatorname{coker}(\mathrm{loc}/L)
       =\dim Q-\dim\Sha^1_S(\mathbf Q,M^D).
                                                               \tag{CM.17}
\]
For clarity, the annihilator of the image of global
$H^1(M)$ in the local $H^1(M)$ is the image of global
$H^1(M^D)$.  Intersecting this image with $L^\perp$
identifies the dual of the cokernel on the left with
$Q/\Sha^1_S(M^D)$.  Also
$\dim\Sha^2_S(M)=\dim\Sha^1_S(M^D)$ by Poitou--Tate.
Taking dimensions in (CM.15) therefore gives
$\dim\mathcal O_\psi=\dim Q$.

The group $Q$ is strict at $p$ and unramified at every
other finite prime.  Restriction to $G_F$ and class field
theory identify it with
\[
              Q=\operatorname{Hom}_{\Delta_0}
                          (A_F/pA_F,k(\omega^{1-e})).           \tag{CM.18}
\]
Restriction and descent are exact because $p\nmid|\Delta_0|$.
Conversely, an everywhere-unramified class-group homomorphism
is zero on the full decomposition group at the unique prime
over $p$: its Frobenius is the class of a principal prime.
Local restriction is injective, so the descended class is
strict at $\mathbf Q_p$.  No quotient by the prime classes
over $D$ occurs in (CM.18); only unramifiedness, rather
than zero localization, is prescribed there.
This proves (CM.7).  We assert this exact dimension and the
canonical cochain obstruction (CM.14); we do not obtain a
canonical class-group-valued functional merely by comparing
two short exact sequences.

Finally the carried Herbrand implication, explicitly (AB.12)
in KT-RATIONAL-POSITIVE-ROOT-UNIT, kills
$A_F^{\omega^{1-e}}$ when $B_b/b$ is a unit.
Thus (CM.7) makes $\mathcal O_\psi$ zero, proving the last
assertion.
\end{proof}

\paragraph{An explicit repeated Massey identity.}
The rational obstruction is not just assigned a character.
With pointwise product $uv$ define the two one-cochains
\[
 a_{12}=-w-\tfrac12uv,\qquad
 a_{23}=w-\tfrac12uv.
                                                               \tag{CM.19}
\]
Equation (CM.8) and the cocycle laws give
$da_{12}=u\smile v$ and $da_{23}=v\smile u$.
Thus
\[
 \Theta(u,v,u)=u\smile a_{23}+a_{12}\smile u
                      \in Z^2(G_{\mathbf Q,S},k(\psi))
                                                               \tag{CM.20}
\]
is a specified defining value of the repeated triple
Massey product.  Direct expansion gives the exact cochain
identity
\[
             \Omega_+=-\tfrac32\Theta(u,v,u)
                               +\tfrac14d(u^2v).              \tag{CM.21}
\]
Consequently the raw image of (CM.6) is
$\tfrac32[\Theta(u,v,u)]$.
The usual unrestricted defining-system indeterminacy is not
silently discarded: the ordinary quadratic normalization
above fixes the allowed $w$ modulo global coboundary.
For fixed $[u]$, the resulting relative obstruction is linear
in $[v]$, and scaling $[u]$ by $\lambda$ scales it by
$\lambda^2$.  This follows either from (CM.8)--(CM.10) by
choosing linear cochain splittings, whose choices do not
affect (CM.6), or from the two-parameter grading.

\paragraph{What would be needed for an exclusion.}
The first surviving arithmetic requirement is therefore
the vanishing of this particular relative cubic class,
whose dimension is the rational $\omega^{1-e}$ class
component.  It is more precise than the existence of a
raw $H^2$ group or an unrestricted Massey defining system.
The three auxiliary units do not evaluate this class.
Nor does a nonzero rational class component prove that
the value (CM.6) is nonzero.

An exclusion would require an actual arithmetic
nonvanishing or injectivity assertion for this map,
together with a proof that the relevant source is forced
to admit the specified two-parameter first-order jet.
Neither assertion is proved here.  In particular this
calculation is not a new upper bound for the original
strict Selmer group, and a characteristic-zero source
point is not inferred from a cubic lift.

\paragraph{Exact verification and imported scope.}
The separate certificate
\texttt{cubic\_massey\_check.py} verifies the finite BCH
identity, both two-cocycle identities, all five involution
signs, the determinant, both Massey defining equations,
and (CM.21) by exact symbolic polynomial arithmetic.
The arithmetic cohomology input (CM.12) is the separately
whole-reviewed KT-CUBIC-TWISTED-COHOMOLOGY.
For (CM.15)--(CM.18) the imported duality is Milne,
\emph{Arithmetic Duality Theorems}, second edition,
Chapter~I, Theorem~4.10(a)--(b) and its complete exact
sequence on pp.~56--57.  Local unramified orthogonality
is Theorem~2.6 on p.~30, with its complete proof continuing
to p.~31; at $q\ne p$ the finite dual takes values in
unramified units, so its $M^d$ equals $M^D$ here.
Those complete passages were read for this application.
The foundational duality theorems are imported.
